%% kapa: exam=1 round=36 year=2025 period=1 form=A subject=civil_law subject_ko=민법 seq=16 number=16
%% source: Q-Net 감정평가사 기출문제 (cst00309) artlSeq=5231525 file=raw/1008/5231525/2025년 제36회 감정평가사 1차시험 1교시 문제지 원본.pdf
%% conversion: kordoc 4.24.51 -> tools/gen1.py (auto); review=proofread-sonnet; flags=none
\begin{question}{16}
소멸시효의 중단에 관한 설명으로 옳지 않은 것은? (다툼이 있으면 판례에 따름)
\begin{choices}
  \item 시효중단의 효력이 미치는 당사자란 중단행위에 관여한 당사자를 말한다.
  \item 원인채권의 지급을 확보하기 위한 방법으로 어음이 수수된 경우, 원인채권에 기하여 청구한 것만으로는 어음채권의 소멸시효를 중단시키지 못한다.
  \item 가압류의 피보전채권에 관하여 본안의 승소판결이 확정되었다면 가압류에 의한 시효중단의 효력은 재판상 청구에 흡수되어 소멸된다.
  \item 요역지가 수인의 공유인 경우에 그 1인에 의한 지역권소멸시효의 중단은 다른 공유자를 위하여 효력이 있다.
  \item 소멸시효의 중단사유로서의 승인은 소멸시효의 진행이 개시된 이후에만 가능하다.
\end{choices}
\end{question}
\answer{3}
